% Source slide 102
\section[4.7 Change of Basis]{4.7 Change of Basis}
\chaptercard{4}{Determinants}{4.7}{CHANGE OF BASIS}

% Source slide 103
\begin{frame}[shrink=8]{CHANGE OF BASIS}
\begin{itemize}
\item \textbf{Example 1} Consider two bases $\beta=\{\vect b_1,\vect b_2\}$ and $\mathcal C=\{\vect c_1,\vect c_2\}$ for a vector space $V$, such that
\begin{equation*}
\vect b_1=4\vect c_1+\vect c_2\quad\text{and}\quad
\vect b_2=-6\vect c_1+\vect c_2. \tag{1}
\end{equation*}
\item Suppose
\begin{equation*}
\vect x=3\vect b_1+\vect b_2. \tag{2}
\end{equation*}
\item That is, suppose $[\vect x]_{\beta}=\begin{bmatrix}3\\1\end{bmatrix}$. Find $[\vect x]_{\mathcal C}$.
\end{itemize}
\end{frame}

% Source slide 104
\begin{frame}[shrink=8]{CHANGE OF BASIS}
\begin{itemize}
\item \textbf{Solution} Apply the coordinate mapping determined by $\mathcal C$ to $\vect x$ in (2). Since the coordinate mapping is a linear transformation,
\begin{align*}
[\vect x]_{\mathcal C}&=[3\vect b_1+\vect b_2]_{\mathcal C}\\
&=3[\vect b_1]_{\mathcal C}+[\vect b_2]_{\mathcal C}.
\end{align*}
\item We can write the vector equation as a matrix equation, using the vectors in the linear combination as the columns of a matrix:
\begin{equation*}
[\vect x]_{\mathcal C}=
\begin{bmatrix}[\vect b_1]_{\mathcal C}&[\vect b_2]_{\mathcal C}\end{bmatrix}
\begin{bmatrix}3\\1\end{bmatrix}. \tag{3}
\end{equation*}
\end{itemize}
\end{frame}

% Source slide 105
\begin{frame}[shrink=8]{CHANGE OF BASIS}
\begin{itemize}
\item This formula gives $[\vect x]_{\mathcal C}$, once we know the columns of the matrix. From (1),
\[
[\vect b_1]_{\mathcal C}=\begin{bmatrix}4\\1\end{bmatrix}
\quad\text{and}\quad
[\vect b_2]_{\mathcal C}=\begin{bmatrix}-6\\1\end{bmatrix}.
\]
\item Thus, (3) provides the solution:
\[
[\vect x]_{\mathcal C}=
\begin{bmatrix}4&-6\\1&1\end{bmatrix}
\begin{bmatrix}3\\1\end{bmatrix}
=\begin{bmatrix}6\\4\end{bmatrix}.
\]
\item The $\mathcal C$-coordinates of $\vect x$ match those of the $\vect x$ in Fig. 1, as seen on the next slide.
\end{itemize}
\end{frame}

% Source slide 106
\begin{frame}{CHANGE OF BASIS}
\centering
\begin{tikzpicture}[x=0.68cm,y=0.68cm,>=Stealth,font=\scriptsize]
  \begin{scope}[shift={(-7.0,0)}]
    \clip (-3.1,-2.2) rectangle (3.1,2.2);
    \draw[thin] (-3.1,-2.2) rectangle (3.1,2.2);
    \coordinate (O) at (-1.55,0);
    \coordinate (b1) at (1.10,-0.45);
    \coordinate (b2) at (-0.65,1.10);
    \foreach \i in {-3,...,5} {\draw[gray] ($(O)+\i*(b1)-3*(b2)$)--($(O)+\i*(b1)+3*(b2)$);}
    \foreach \j in {-3,...,4} {\draw[gray] ($(O)-4*(b1)+\j*(b2)$)--($(O)+5*(b1)+\j*(b2)$);}
    \draw[diagramblue,thick] ($(O)-2*(b1)$)--($(O)+4*(b1)$);
    \draw[diagramblue,thick] ($(O)-2*(b2)$)--($(O)+3*(b2)$);
    \fill[diagramblue!25] (O)--($(O)+3*(b1)$)--($(O)+3*(b1)+(b2)$)--($(O)+(b2)$)--cycle;
    \draw[->,diagramblue,very thick] (O)--($(O)+(b1)$) node[midway,below] {$\vect b_1$};
    \draw[->,diagramblue,very thick] (O)--($(O)+(b2)$) node[midway,left] {$\vect b_2$};
    \fill[diagramblue] ($(O)+3*(b1)+(b2)$) circle (2pt) node[above right] {$\vect x$};
    \node[below left] at (O) {$0$};
    \node[below] at ($(O)+3*(b1)$) {$3\vect b_1$};
  \end{scope}
  \begin{scope}[shift={(3.1,0)}]
    \clip (-3.1,-2.2) rectangle (3.1,2.2);
    \draw[thin] (-3.1,-2.2) rectangle (3.1,2.2);
    \coordinate (P) at (-1.60,0);
    \coordinate (c1) at (0.28,-0.30);
    \coordinate (c2) at (0.42,0.28);
    \foreach \i in {-5,...,9} {\draw[gray] ($(P)+\i*(c1)-5*(c2)$)--($(P)+\i*(c1)+6*(c2)$);}
    \foreach \j in {-6,...,7} {\draw[gray] ($(P)-6*(c1)+\j*(c2)$)--($(P)+9*(c1)+\j*(c2)$);}
    \draw[diagramblue,thick] ($(P)-3*(c1)$)--($(P)+8*(c1)$);
    \draw[diagramblue,thick] ($(P)-4*(c2)$)--($(P)+6*(c2)$);
    \fill[diagramblue!25] (P)--($(P)+6*(c1)$)--($(P)+6*(c1)+4*(c2)$)--($(P)+4*(c2)$)--cycle;
    \draw[->,diagramblue,very thick] (P)--($(P)+(c1)$) node[midway,below right] {$\vect c_1$};
    \draw[->,diagramblue,very thick] (P)--($(P)+(c2)$) node[midway,above left] {$\vect c_2$};
    \fill[diagramblue] ($(P)+6*(c1)+4*(c2)$) circle (2pt) node[above right] {$\vect x$};
    \node[below left] at (P) {$0$};
    \node[above right] at ($(P)+6*(c1)$) {$6\vect c_1$};
    \node[above] at ($(P)+4*(c2)$) {$4\vect c_2$};
  \end{scope}
  \node at (-7,-2.55) {(a)};
  \node at (3.1,-2.55) {(b)};
\end{tikzpicture}

\vspace{0.25em}
\begin{flushleft}
\textbf{FIGURE 1}\quad Two coordinate systems for the same vector space.
\end{flushleft}
\end{frame}

% Source slide 107
\begin{frame}[shrink=8]{CHANGE OF BASIS}
\begin{itemize}
\item \sourcehead{Theorem 15:} Let $\beta=\{\vect b_1,\ldots,\vect b_n\}$ and $\mathcal C=\{\vect c_1,\ldots,\vect c_2\}$ for a vector space $V$. Then there is a unique $n\times n$ matrix $P_{\mathcal C\leftarrow\beta}$ such that
\begin{equation*}
[\vect x]_{\mathcal C}=P_{\mathcal C\leftarrow\beta}[\vect x]_{\beta}. \tag{4}
\end{equation*}
\item The columns of $P_{\mathcal C\leftarrow\beta}$ are the $\mathcal C$-coordinate vectors of the vectors in the basis $\beta$. That is,
\begin{equation*}
P_{\mathcal C\leftarrow\beta}=
\begin{bmatrix}[\vect b_1]_{\mathcal C}&[\vect b_2]_{\mathcal C}&\cdots&[\vect b_n]_{\mathcal C}\end{bmatrix}. \tag{5}
\end{equation*}
\end{itemize}
\end{frame}

% Source slide 108
\begin{frame}[shrink=5]{CHANGE OF BASIS}
\begin{itemize}
\item The matrix $P_{\mathcal C\leftarrow\beta}$ in Theorem 15 is called the \textbf{change-of-coordinates matrix from $\beta$ to $\mathcal C$}. Multiplication by $P_{\mathcal C\leftarrow\beta}$ converts $\beta$-coordinates into $\mathcal C$-coordinates.
\item Figure 2 below illustrates the change-of-coordinates equation (4).
\end{itemize}
\centering
\begin{tikzpicture}[>=Stealth,font=\scriptsize,scale=.82]
  \fill[diagramblue!22] (-1.5,2.5)--(1.6,2.5)--(2.15,3.8)--(-.95,3.8)--cycle;
  \node at (1.75,3.55) {$V$};
  \fill[diagramblue] (-.65,2.85) circle (2pt) node[above right] {$\vect x$};
  \fill[diagramblue!22] (-5,-.2)--(-1.5,-.2)--(-.9,1.15)--(-4.4,1.15)--cycle;
  \fill[diagramblue!22] (2.15,-.2)--(5.65,-.2)--(6.25,1.15)--(2.75,1.15)--cycle;
  \fill[diagramblue] (-2.15,.72) circle (2pt) node[right] {$[\vect x]_{\mathcal C}$};
  \fill[diagramblue] (4.55,.72) circle (2pt) node[right] {$[\vect x]_{\beta}$};
  \node at (-4.25,.12) {$\R^n$};
  \node at (3.0,.12) {$\R^n$};
  \draw[->,thick] (-.8,2.55)--(-3.15,1.2) node[midway,left] {$[\ ]_{\mathcal C}$};
  \draw[->,thick] (.85,2.55)--(3.85,1.2) node[midway,right] {$[\ ]_{\beta}$};
  \draw[->,thick] (3.1,.45)--(-1.6,.45) node[midway,above,align=center] {multiplication\\by $P_{\mathcal C\leftarrow\beta}$};
\end{tikzpicture}

\begin{flushleft}
\textbf{FIGURE 2}\quad Two coordinate systems for $V$.
\end{flushleft}
\end{frame}

% Source slide 109
\begin{frame}[shrink=9]{CHANGE OF BASIS}
\begin{itemize}
\item The columns of $P_{\mathcal C\leftarrow\beta}$ are linearly independent because they are the coordinate vectors of the linearly independent set $\beta$.
\item Since $P_{\mathcal C\leftarrow\beta}$ is square, it must be invertible, by the Invertible Matrix Theorem. Left-multiplying both sides of equation (4) by $(P_{\mathcal C\leftarrow\beta})^{-1}$ yields
\[
(P_{\mathcal C\leftarrow\beta})^{-1}[\vect x]_{\mathcal C}=[\vect x]_{\beta}.
\]
\item Thus $(P_{\mathcal C\leftarrow\beta})^{-1}$ is the matrix that converts $\mathcal C$-coordinates into $\beta$--coordinates. That is,
\begin{equation*}
(P_{\mathcal C\leftarrow\beta})^{-1}=P_{\beta\leftarrow\mathcal C}. \tag{6}
\end{equation*}
\end{itemize}
\end{frame}

% Source slide 110
\begin{frame}[shrink=10]{CHANGE OF BASIS IN $\R^n$}
\begin{itemize}
\item If $\beta=\{\vect b_1,\ldots,\vect b_n\}$ and $\mathcal E$ is the standard basis $\{\vect e_1,\ldots,\vect e_n\}$ in $\R^n$, then $[\vect b_1]_{\mathcal E}=\vect b_1$, and likewise for the other vectors in $\beta$. In this case, $P_{\mathcal E\leftarrow\beta}$ is thesame as the change-of-coordinates matrix $P_{\beta}$ introduced in Section 4.4, namely,
\[
P_{\beta}=\begin{bmatrix}\vect b_1&\vect b_2&\cdots&\vect b_n\end{bmatrix}.
\]
\item To change coordinates between two nonstandard bases in $\R^n$,we need Theorem 15. The theorem shows that to solve the change-of-basis problem, we need the coordinate vectors of the old basis relative to the new basis.
\end{itemize}
\end{frame}

% Source slide 111
\begin{frame}[shrink=12]{CHANGE OF BASIS IN $\R^n$}
\begin{itemize}
\item \textbf{Example 2} Let
$\vect b_1=\begin{bmatrix}-9\\1\end{bmatrix}$,
$\vect b_2=\begin{bmatrix}-5\\-1\end{bmatrix}$,
$\vect c_1=\begin{bmatrix}1\\-4\end{bmatrix}$,
$\vect c_2=\begin{bmatrix}3\\-5\end{bmatrix}$
and consider the bases for $\R^n$ given by $\beta=\{\vect b_1,\vect b_2\}$ and $\mathcal C=\{\vect c_1,\vect c_2\}$. Find the change-of-coordinates matrix from $\beta$ to $\mathcal C$.
\item \textbf{Solution} The matrix $P_{\beta\leftarrow\mathcal C}$ involves the $\mathcal C$-coordinate vectors of $\vect b_1$ and $\vect b_2$. Let $[\vect b_1]_{\mathcal C}=\begin{bmatrix}x_1\\x_2\end{bmatrix}$ and $[\vect b_2]_{\mathcal C}=\begin{bmatrix}y_1\\y_2\end{bmatrix}$. Then, by definition,
\[
\begin{bmatrix}\vect c_1&\vect c_2\end{bmatrix}\begin{bmatrix}x_1\\x_2\end{bmatrix}=\vect b_1
\quad\text{and}\quad
\begin{bmatrix}\vect c_1&\vect c_2\end{bmatrix}\begin{bmatrix}y_1\\y_2\end{bmatrix}=\vect b_2.
\]
\end{itemize}
\end{frame}

% Source slide 112
\begin{frame}[shrink=10]{CHANGE OF BASIS IN $\R^n$}
\begin{itemize}
\item To solve both systems simultaneously, augment the coefficient matrix with $\vect b_1$ and $\vect b_2$, and row reduce:
\begin{equation*}
\left[\begin{array}{cc|cc}\vect c_1&\vect c_2&\vect b_1&\vect b_2\end{array}\right]=
\left[\begin{array}{cc|cc}1&3&-9&-5\\-4&-5&1&-1\end{array}\right]
\sim
\left[\begin{array}{cc|cc}1&0&6&4\\0&1&-5&-3\end{array}\right]. \tag{7}
\end{equation*}
\item Thus
\[
[\vect b_1]_{\mathcal C}=\begin{bmatrix}6\\-5\end{bmatrix}
\quad\text{and}\quad
[\vect b_2]_{\mathcal C}=\begin{bmatrix}4\\-3\end{bmatrix}.
\]
\item The desired change-of-coordinates matrix is therefore
\[
P_{\mathcal C\leftarrow\beta}=
\begin{bmatrix}[\vect b_1]_{\mathcal C}&[\vect b_2]_{\mathcal C}\end{bmatrix}
=\begin{bmatrix}6&4\\-5&-3\end{bmatrix}.
\]
\end{itemize}
\end{frame}
